% Compile with: uplatex file.tex && dvipdfmx file.dvi
%           or: latexmk -uplatex file.tex
\documentclass[dvipdfmx]{article}
\usepackage{amsmath, amssymb}
\usepackage{tikz}

\title{scenario\_transverse\_field\_ising}
\date{\today}

\begin{document}
\maketitle

Consider the two-site transverse-field Ising model. Prepare the coupling constant $J$ (the strength of the Ising interaction) and the transverse-field strength $\Gamma$ as scalar parameters.

Define J, $\Gamma$ as scalars

\[
  J, \Gamma
\]

Prepare the spin operator $S$ used on both sites.

S as the spin operator

\[
  S, S^{x}, S^{y}, S^{z}
\]

Register the transverse-field Ising Hamiltonian $H=-JS^z_1S^z_2-\Gamma(S^x_1+S^x_2)$ (the Ising coupling term and the transverse-field term).

$H = -J Sz_1 Sz_2 - \Gamma (Sx_1 + Sx_2)$

\begin{equation}
  H = -J Sz_1 Sz_2 - \Gamma (Sx_1 + Sx_2)
  \label{eq:E1}
\end{equation}

Compute the commutator of the Hamiltonian $H$ and $S^z$ on site 1, to check part of the equation of motion following from the SU(2) commutation relations.

[E1, $Sz_1$]

\begin{equation}
  2 \Gamma \mathit{i} Sy_1
  \label{eq:E2}
\end{equation}

To carry out a mean-field approximation, prepare the mean field (a c-number) $m_z=\langle S^z\rangle$, which will replace the spin on site 2, as a scalar parameter.

Define mz as a scalar

\[
  mz
\]

Replace only the operator $S^z_2$ on site 2 with the mean field $m_z$, obtaining the mean-field Hamiltonian $H_{\rm MF}$ (the operators on site 1 and the transverse-field term should be unchanged).

In E1, substitute mz for $S^z_2$

\begin{equation}
  H = - \Gamma Sx_1 - \Gamma Sx_2 - J mz Sz_1
  \label{eq:E3}
\end{equation}

\end{document}